\documentclass[11pt, letterpaper]{scrartcl}
\usepackage{azzam}

\title{Campur}

\begin{document}
\begin{enumerate} 
    \item Tentukan semua fungsi \( f: \mathbb{N} \to \mathbb{N} \) dengan sifat
       \[
       (f(m))^2 + f(n) \text{ habis membagi } (m^2 + n)^2
       \]
       untuk semua \( m, n \in \mathbb{N} \).
    
    \item Papan catur \( 8 \times 8 \) diwarnai hitam putih seperti biasa. Tentukan banyaknya cara meletakkan 16 koin identik pada kotak-kotak putih sehingga setiap kotak berisi paling banyak satu koin dan tidak ada dua koin yang bersebelahan secara diagonal.
    
    \item Sebuah segitujuh beraturan tiap titik sudutnya diwarnai merah atau biru. Buktikan bahwa dapat dipilih tiga titik sudut dengan warna yang sama sehingga membentuk segitiga samakaki.

    \item Buktikan bahwa tidak terdapat pasangan bilangan bulat positif $(a, b)$ sehingga $a + b$, $a + 2b$, dan $2a + b$ ketiganya kuadrat sempurna.
    
    \item Sebuah polinomial $P(x)$ memenuhi $P(a) = P(b) = P(c) = 1$ untuk suatu bilangan bulat berbeda $a$, $b$, dan $c$. Buktikan bahwa $P(x)$ tidak mempunyai akar bilangan bulat.
    
    \item Tentukan semua fungsi kontinu yang memenuhi
        $$x^2f(y) + y^2f(x) = (x + y)f(x)f(y)$$
        untuk setiap bilangan real $x$ dan $y$.
        
    \item Diberikan segitiga samakaki $ABC$ dengan $AB = BC$. Titik $D$ dan $E$ berturut-turut pada sisi $AB$ dan $BC$ sehingga $AD + CE = DE$. Sebuah garis sejajar $BC$ melewati titik tengah segmen $DE$ memotong $AC$ di $F$. Tunjukkan bahwa $\angle DFE = 90^\circ$.

    % https://www.cut-the-knot.org/arithmetic/algebra/AsymmetricInequality.shtml#proof
    \item Untuk sebarang bilangan riil positif $x,y,z$, buktikan bahwa
    \begin{align*}
        (x+y)(x+z) \ge \sqrt{4xyz(x+y+z)}
    \end{align*}

    % paket soal OSP : NT
    \item Misalkan tripel bilangan prima $(p,q,r)$ memenuhi $3p^4-5q^4-4r^2=26$. Jika $S = p+q+r$, hitunglah jumlah semua $S$ yang mungkin.

    % paket soal OSP : kombin
    \item Terdapat 20 anggota di suatu klub tenis yang menjadwalkan tepat 14 permainan antar dua orang diantara mereka dengan setiap anggota klub bermain minimal satu kali. Buktikan bahwa dalam pembagian ini, terdapat himpunan 6 permainan dengan 12 pemain yang berbeda.

\end{enumerate}
\end{document}